% ============================================================================
%  Anexos: código fuente
%  El código se lee directamente de ../firmware/, así que cualquier cambio en
%  el firmware aparece solo en el informe.
% ============================================================================
\section{Firmware principal}
\label{anx:codigo}

Archivo \texttt{firmware/taller2\_robot\_diferencial.c} (configuración
\texttt{completo}). Se compila con MPLAB XC8 para el PIC16F877A.

\lstinputlisting[caption={Firmware del robot móvil diferencial
(\texttt{taller2\_robot\_diferencial.c}).}, label={lst:firmware}]
{../firmware/taller2_robot_diferencial.c}

\section{Código de la guía del taller}
\label{anx:docente}

Archivo \texttt{firmware/codigo\_docente.c} (configuración \texttt{try1}): el
código entregado por el docente, sin los encabezados de página que se
copiaron desde el PDF.

\lstinputlisting[caption={Código de la guía (\texttt{codigo\_docente.c}).},
label={lst:docente}]{../firmware/codigo_docente.c}

\section{Prueba de montaje con LEDs}
\label{anx:leds}

Archivo \texttt{firmware/prueba\_leds.c} (configuración \texttt{leds}).

\lstinputlisting[caption={Prueba de montaje (\texttt{prueba\_leds.c}).},
label={lst:leds}]{../firmware/prueba_leds.c}
